\section{The Problem}
\label{sec:problem}

\subsection{Changing the order of axes is free}

Deep learning code changes the order of a tensor's dimensions often. In
PyTorch, \texttt{permute} and \texttt{transpose} do
this~\cite{pytorch_permute_docs, pytorch_transpose_docs}. They are free. They
change a label. They do not move any data.

\subsection{Fixing the layout is not free}

After that, the data is in the wrong order in memory. Some operations will not
run on it. Others run, but they use a slow path. They do not warn you.

The normal fix is \texttt{.contiguous()}~\cite{pytorch_contiguous_2016}. It
makes a copy. It asks for a new tensor. It writes every value into the right
place. Then it throws the old tensor away.

\subsection{The copy is the problem}

While the copy runs, both tensors are in memory. So a tensor of $B$ bytes needs
at least $2B$ bytes. The old one must stay until it has been read. The new one
must exist to be written into. This is what a copy means. You cannot avoid it
by changing the loop.

If the machine is almost full, the job fails here. A 2026 PyTorch design
document says non-contiguous tensors now appear all over modern
training~\cite{pytorch_rfc_comms_2026}. It says some communication routines
will not take them. It says the usual fixes either make a full copy or need a
buffer as big as the data.

There is a second problem. The copy is a large block of memory. It only lives
for a short time. It appears in the middle of work that already uses a lot of
memory. Earlier work shows this pattern causes
fragmentation~\cite{rajbhandari2020zero}. Memory is free, but it is in small
pieces. The next request needs one big piece and fails.

\subsection{What we do}

We fix the layout in place. We never ask for a second tensor. We break the
reordering into small swaps. Each swap exchanges two groups of axes that sit
next to each other. We store only a few hundred bytes while we work.

Doing this in place is not new. Section~\ref{sec:related} lists the earlier
work. Two things here are new. In more than two dimensions, this can do
\emph{less work than a copy}. And we can tell which cases those are by looking
at the shape, before we touch any data.
